\chapter{Thyrotoxicosis}
\index{Thyrotoxicosis}

\medskip
\noindent\textit{Educational reference; always seek professional advice for individual care.}

\section*{Definition and Overview}
Thyrotoxicosis is the clinical state caused by excess thyroid hormone action in tissues, regardless of whether the thyroid gland itself is overproducing hormone. Hyperthyroidism is the subset in which the gland synthesises too much hormone (for example Graves disease or toxic nodules). Thyrotoxicosis also includes release of stored hormone in thyroiditis and excess intake of thyroid tablets (factitious or iatrogenic). People experience a hypermetabolic syndrome: heat intolerance, weight loss, tremor, palpitations, anxiety, and diarrhoea, with serious risks of atrial fibrillation, bone loss, and, rarely, thyroid storm. Accurate cause-finding guides therapy because antithyroid drugs, radioiodine, surgery, and simple supportive care are not interchangeable. This chapter covers recognition, investigations, treatment options, and emergency features relevant to local care.

\section*{Epidemiology}
Thyrotoxicosis is commoner in women. Graves disease predominates in younger and middle-aged adults; toxic multinodular goitre and toxic adenoma become more frequent with age and account for many older presentations, including new atrial fibrillation. Transient thyrotoxicosis from thyroiditis is under-recognised when pain is absent. Iodine loads from contrast media or amiodarone precipitate disease in susceptible glands. In pregnancy, differential diagnosis includes gestational transient thyrotoxicosis and Graves disease, with distinct risks. Subclinical hyperthyroidism (low TSH, normal free hormones) is prevalent in older populations with nodules and carries atrial fibrillation and osteoporosis concerns even without florid symptoms.

\section*{Aetiology and Pathophysiology}
\begin{itemize}
    \item \textbf{Graves disease:} TSH-receptor stimulating antibodies raise hormone output and can drive orbitopathy and pretibial myxoedema
    \item \textbf{Toxic multinodular goitre and toxic adenoma:} somatic activating mutations create TSH-independent tissue
    \item \textbf{Thyroiditis leak:} follicular damage dumps T4/T3; synthesis is not increased; uptake is low
    \item \textbf{Excess exogenous hormone:} over-replacement or misuse of levothyroxine/T3
    \item \textbf{TSH-secreting pituitary adenoma (rare):} high hormones with inappropriately non-suppressed TSH
    \item \textbf{hCG-driven:} marked hyperemesis or trophoblastic disease
    \item \textbf{Iodine-induced (Jod--Basedow):} autonomous tissue overproduces when iodine-repleted
    \item \textbf{Tissue effects:} increased basal metabolic rate, heightened catecholamine sensitivity, shortened cardiac refractory periods, increased bone turnover
\end{itemize}

\section*{Clinical Features}
\begin{itemize}
    \item \textbf{Constitutional:} weight loss despite appetite, heat intolerance, sweating, fatigue from insomnia and weakness
    \item \textbf{Neuropsychiatric:} irritability, anxiety, poor concentration; apathetic thyrotoxicosis in some elderly
    \item \textbf{Cardiovascular:} sinus tachycardia, palpitations, atrial fibrillation, exertional dyspnoea, high-output strain
    \item \textbf{Neuromuscular:} fine tremor, proximal myopathy, hyper-reflexia; rare periodic paralysis in susceptible groups
    \item \textbf{Gastrointestinal:} frequency of stool, thirst
    \item \textbf{Reproductive:} oligomenorrhoea, reduced fertility; gynaecomastia occasionally in men
    \item \textbf{Neck:} diffuse goitre (Graves), nodular goitre, or tender gland (subacute thyroiditis)
    \item \textbf{Eyes:} lid lag in any thyrotoxicosis; true orbitopathy points to Graves
    \item \textbf{Thyroid storm features:} fever, delirium, vomiting, diarrhoea, tachycardia out of proportion---emergency
\end{itemize}

\section*{Differential Diagnosis}
\begin{itemize}
    \item \textbf{Primary anxiety and panic disorder}
    \item \textbf{Phaeochromocytoma and cocaine or amphetamine use}
    \item \textbf{Menopause and anaemia}
    \item \textbf{Malignancy-related weight loss}
    \item \textbf{Chronic infection}
    \item \textbf{Graves versus thyroiditis versus toxic nodule:} critical therapeutic fork
    \item \textbf{Exogenous hormone versus endogenous production:} thyroglobulin low when hormone is swallowed
\end{itemize}

\section*{When to Seek Medical Care}
See a GP promptly for persistent palpitations, unexplained weight loss, tremor, or new AF. Urgent same-day care is needed for marked agitation, very high heart rate, or pregnancy with thyrotoxic symptoms.

\textbf{Contact your local emergency services for:}
\begin{itemize}
    \item Chest pain, severe breathlessness, fainting, or stroke-like symptoms with known or suspected thyrotoxicosis
    \item High fever, confusion, vomiting, and extreme tachycardia suggesting thyroid storm
    \item Sudden severe eye pain or vision loss in Graves disease
\end{itemize}

\section*{Diagnosis and Investigations}
\begin{itemize}
    \item \textbf{TSH suppressed with high free T4 and/or free T3:} biochemical confirmation of thyrotoxicosis
    \item \textbf{TRAb:} supports Graves disease
    \item \textbf{ESR/CRP and clinical tenderness:} subacute thyroiditis clues
    \item \textbf{Uptake scan (when not pregnant/breastfeeding):} high diffuse (Graves), focal (toxic nodule), low (thyroiditis or exogenous hormone)
    \item \textbf{Ultrasound:} vascularity and nodularity; first-line structural test in pregnancy instead of nuclear scans
    \item \textbf{ECG:} rhythm assessment
    \item \textbf{FBC and liver enzymes:} baseline before antithyroid drugs
    \item \textbf{Beta-hCG:} when pregnancy possible
\end{itemize}

\section*{Management}
\begin{itemize}
    \item \textbf{Beta-blockers:} rapid symptom control (for example propranolol) unless contraindicated
    \item \textbf{Antithyroid drugs:} carbimazole first-line for many Graves cases in many countries; PTU preferred in first trimester pregnancy when drugs are needed; educate about agranulocytosis and liver toxicity
    \item \textbf{Radioiodine:} definitive therapy for Graves or toxic nodules after suitability assessment; avoid in pregnancy and many with active significant orbitopathy without specialist precautions
    \item \textbf{Thyroidectomy:} large goitres, suspected cancer, pregnancy second-trimester selected cases, or patient preference after rendering euthyroid
    \item \textbf{Thyroiditis:} supportive care and beta-blockade; steroids for severe subacute pain; no ATD for pure leak
    \item \textbf{Exogenous excess:} dose reduction
    \item \textbf{Thyroid storm:} ICU-level care with ATDs, iodine after ATDs, steroids, beta-blockade, supportive cooling, and treatment of triggers
    \item \textbf{AF management:} rate/rhythm and anticoagulation decisions per cardiology guidance while rendering euthyroid
    \item \textbf{Bone and weight recovery:} nutrition and falls prevention in older adults
\end{itemize}

\section*{Complications}
\begin{itemize}
    \item \textbf{Atrial fibrillation and thromboembolism}
    \item \textbf{High-output heart failure}
    \item \textbf{Osteoporosis and fractures}
    \item \textbf{Thyroid storm mortality:} still significant without aggressive care
    \item \textbf{Antithyroid drug agranulocytosis and hepatotoxicity}
    \item \textbf{Radioiodine-induced hypothyroidism:} expected often; needs replacement
    \item \textbf{Worsening orbitopathy:} especially in smokers after radioiodine without cover
    \item \textbf{Pregnancy loss and fetal hyperthyroidism:} in uncontrolled or antibody-mediated disease
\end{itemize}

\section*{Prognosis and Outcomes}
With timely treatment, most people return to normal metabolic health. Graves may remit after 12--18 months of drugs in a substantial minority; others need definitive therapy. Toxic nodules rarely remit without ablative treatment. Thyroiditis-related thyrotoxicosis is self-limited but may flip to hypothyroidism. Older adults with delayed recognition face more cardiac morbidity. Adherence to monitoring after any definitive therapy underpins long-term success.

\section*{Prevention and Self-Care}
\begin{itemize}
    \item Do not adjust thyroid tablet doses without blood tests and advice
    \item Take antithyroid drugs exactly as prescribed; stop and seek urgent FBC if fever and sore throat occur
    \item Stop smoking; protect eyes with lubricants if advised
    \item Limit excess iodine supplements unless prescribed
    \item Attend early pregnancy care if you have a history of Graves disease---antibodies can persist
    \item Seek AF review rather than attributing all palpitations to anxiety
    \item Carry an updated medicine list for emergency presentations
\end{itemize}

\section*{Sources and Further Reading}
\begin{itemize}
    \item NHS. \textit{Overactive thyroid (hyperthyroidism).} https://www.nhs.uk/
    \item Mayo Clinic. \textit{Hyperthyroidism.} https://www.mayoclinic.org/
    \item MSD Manual. \textit{Hyperthyroidism.} https://www.msdmanuals.com/
    \item MedlinePlus. \textit{Hyperthyroidism.} https://medlineplus.gov/
    \item Better Health Channel. \textit{Hyperthyroidism.} https://www.betterhealth.vic.gov.au/
    \item Healthdirect. \textit{Hyperthyroidism.} https://www.healthdirect.gov.au/
\end{itemize}
